% ============================================================
%  PEMBAHASAN SUPER DETAIL — LATIHAN SOAL & STUDI KASUS
%  BAB 7: STATISTIKA
%  Catatan: Soal Akhir Bab TIDAK dibahas di sini (sesuai permintaan).
% ============================================================
% ============================================================
%  PREAMBUL BERSAMA — BUKU PEMBAHASAN LATIHAN SOAL
%  Dipakai oleh MAE-2026-2027-Latihan1.tex ... Latihan9.tex
%  Kompilasi: pdfLaTeX (disarankan Overleaf / MiKTeX)
% ============================================================
\documentclass[11pt,a4paper]{report}

% ---------- PAKET ----------
\usepackage[utf8]{inputenc}
\usepackage[T1]{fontenc}
\usepackage[bahasa]{babel}
\usepackage{amsmath,amssymb,amsthm,mathtools}
\usepackage{geometry}
\usepackage{xcolor}
\usepackage{graphicx}
\usepackage{enumitem}
\usepackage{tikz}
\usetikzlibrary{calc,arrows.meta,angles,quotes,positioning,decorations.pathreplacing,patterns}
\usepackage{pgfplots}
\pgfplotsset{compat=1.18}
\usepackage[most]{tcolorbox}
\usepackage{fancyhdr}
\usepackage{titlesec}
\usepackage{array,booktabs}
\usepackage{multicol}
\usepackage{pifont}
\usepackage[colorlinks=true,linkcolor=biru,urlcolor=biru]{hyperref}

\geometry{a4paper,margin=2.5cm,headheight=15pt}

% ---------- WARNA ----------
\definecolor{biru}{HTML}{1F4E79}
\definecolor{birumuda}{HTML}{D6E4F0}
\definecolor{hijau}{HTML}{2E7D32}
\definecolor{hijaumuda}{HTML}{E3F2E1}
\definecolor{oranye}{HTML}{C75B12}
\definecolor{oranyemuda}{HTML}{FBE8DA}
\definecolor{abu}{HTML}{F2F2F2}
\definecolor{ungu}{HTML}{6A1B9A}
\definecolor{ungumuda}{HTML}{EDE2F3}
\definecolor{merah}{HTML}{C62828}
\definecolor{merahmuda}{HTML}{FCE4E4}
\definecolor{tosca}{HTML}{00838F}
\definecolor{toscamuda}{HTML}{E0F2F1}
\definecolor{emas}{HTML}{8D6E00}
\definecolor{emasmuda}{HTML}{FFF6D6}

% ---------- HEADER / FOOTER ----------
\newcommand{\babjudul}{Pembahasan Latihan}
\pagestyle{fancy}
\fancyhf{}
\renewcommand{\headrulewidth}{0.8pt}
\renewcommand{\footrulewidth}{0pt}
\fancyhead[L]{\small\textcolor{ungu}{\leftmark}}
\fancyhead[R]{\small\textcolor{ungu}{\babjudul}}
\fancyfoot[C]{\thepage}

% ---------- GAYA JUDUL BAB ----------
\titleformat{\chapter}[display]
  {\normalfont\huge\bfseries\color{ungu}}
  {\filright\large PEMBAHASAN}{8pt}{\Huge}
\titlespacing*{\chapter}{0pt}{0pt}{20pt}

% ---------- LINGKUNGAN TEOREMA ----------
\newtheoremstyle{gaya}{6pt}{6pt}{}{}{\bfseries\color{biru}}{.}{.5em}{}
\theoremstyle{gaya}
\newtheorem{definisi}{Definisi}[chapter]
\newtheorem{sifat}{Sifat}[chapter]

% ---------- KOTAK: SOAL (ungu) ----------
\newtcolorbox{soalbox}[1]{
  enhanced, breakable, colback=ungumuda, colframe=ungu, boxrule=0.8pt,
  arc=3pt, left=10pt, right=10pt, top=8pt, bottom=8pt,
  title={\textbf{Soal #1}}, fonttitle=\bfseries\color{white}, coltitle=white,
  attach boxed title to top left={xshift=10pt,yshift=-10pt},
  boxed title style={colback=ungu,arc=2pt}}

% ---------- KOTAK: KONSEP KUNCI (hijau) ----------
\newtcolorbox{ingat}[1]{
  enhanced, breakable, colback=hijaumuda, colframe=hijau, boxrule=0.6pt,
  arc=3pt, left=10pt, right=10pt, top=6pt, bottom=6pt,
  title={\textbf{Konsep Kunci: #1}}, fonttitle=\bfseries\color{white}, coltitle=white,
  attach boxed title to top left={xshift=10pt,yshift=-10pt},
  boxed title style={colback=hijau,arc=2pt}}

% ---------- KOTAK: JEBAKAN (merah) ----------
\newtcolorbox{jebakan}{
  enhanced, breakable, colback=merahmuda, colframe=merah, boxrule=0.6pt,
  arc=3pt, left=10pt, right=10pt, top=6pt, bottom=6pt,
  title={\textbf{\ding{55}\ Jebakan \& Kesalahan yang Sering Terjadi}},
  fonttitle=\bfseries\color{white}, coltitle=white,
  attach boxed title to top left={xshift=10pt,yshift=-10pt},
  boxed title style={colback=merah,arc=2pt}}

% ---------- KOTAK: VARIASI SOAL (oranye) ----------
\newtcolorbox{variasi}{
  enhanced, breakable, colback=oranyemuda, colframe=oranye, boxrule=0.6pt,
  arc=3pt, left=10pt, right=10pt, top=6pt, bottom=6pt,
  title={\textbf{\ding{72}\ Bentuk Modifikasi Soal Ini}},
  fonttitle=\bfseries\color{white}, coltitle=white,
  attach boxed title to top left={xshift=10pt,yshift=-10pt},
  boxed title style={colback=oranye,arc=2pt}}

% ---------- KOTAK: STUDI KASUS (biru) ----------
\newtcolorbox{studikasus}[1]{
  enhanced, breakable, colback=abu, colframe=biru, boxrule=0.8pt,
  arc=3pt, left=10pt, right=10pt, top=8pt, bottom=8pt,
  title={\textbf{Studi Kasus: #1}}, fonttitle=\bfseries\color{white}, coltitle=white,
  attach boxed title to top left={xshift=10pt,yshift=-10pt},
  boxed title style={colback=biru,arc=2pt}}

% ---------- KOTAK: TIPS (tosca) ----------
\newtcolorbox{tips}{
  enhanced, breakable, colback=toscamuda, colframe=tosca, boxrule=0.6pt,
  arc=3pt, left=10pt, right=10pt, top=6pt, bottom=6pt,
  borderline west={3pt}{0pt}{tosca}}

% ---------- PERINTAH BANTU ----------
\newcommand{\jawab}{\par\smallskip\noindent\textit{\textcolor{oranye}{Penyelesaian langkah demi langkah:}}\par\smallskip}
\newcommand{\langkah}[1]{\par\smallskip\noindent\textbf{\textcolor{biru}{Langkah #1.}}\ }
\newcommand{\jadi}{\par\smallskip\noindent$\Rightarrow$\ \textbf{Jadi, }}
% Judul tiap nomor soal
\newcommand{\soalno}[2]{%
  \par\bigskip
  \noindent{\large\bfseries\color{ungu}\rule[-2pt]{4pt}{14pt}\ Soal #1}%
  \ {\normalfont\itshape\color{gray}(#2)}\par\nopagebreak\smallskip}

% ---------- HALAMAN JUDUL OTOMATIS ----------
% #1 = judul topik (mis. "Barisan dan Deret"); #2 = nomor bab (angka)
\newcommand{\buathalamanjudul}[2]{%
  \begin{titlepage}
  \centering
  \vspace*{2cm}
  {\color{ungu}\rule{\textwidth}{2pt}}\\[0.6cm]
  {\Huge\bfseries\color{ungu} PEMBAHASAN\\[0.3cm] LATIHAN SOAL}\\[0.4cm]
  {\Large BAB #2 --- #1}\\[0.2cm]
  {\color{ungu}\rule{\textwidth}{2pt}}\\[1.2cm]

  {\large Matematika SMA --- Fase E (Kelas 10)}\\[0.2cm]
  {\large Kurikulum Merdeka}\\[2cm]

  \begin{tcolorbox}[enhanced,width=0.85\textwidth,colback=ungumuda,colframe=ungu,
    arc=4pt,boxrule=1pt]
  \centering\small
  Buku pembahasan ini mengupas \textbf{setiap soal pada kotak ungu (Latihan Soal)}
  di tiap subbab, \emph{selangkah demi selangkah seperti mengajari dari nol}:
  konsep yang dipakai, jebakan yang sering menjerat, cara soal bisa dimodifikasi,
  dan visualisasi. \emph{Soal Akhir Bab tidak dibahas di sini.}
  \end{tcolorbox}

  \vfill
  {\large Tahun Pelajaran 2026/2027}
  \end{titlepage}
  \tableofcontents
  \thispagestyle{empty}
  \clearpage}

% ============================================================
%  PENJAGA KOMPILASI
%  File ini adalah PREAMBUL BERSAMA (di-\input oleh Latihan1..9),
%  BUKAN untuk dikompilasi sendiri. Jika editor terlanjur mengompilasi
%  file ini secara langsung, tampilkan satu halaman info (bukan error).
%  Saat di-\input oleh berkas Latihan, blok ini dilewati.
% ============================================================
\ifnum\pdfstrcmp{\jobname}{MAE-2026-2027-Latihan-preamble}=0
  \begin{document}
  \thispagestyle{empty}
  \begin{center}
  \vspace*{3cm}
  {\Huge\bfseries\color{ungu} Berkas Preambel Bersama}\\[1cm]
  \begin{tcolorbox}[enhanced,width=0.85\textwidth,colback=ungumuda,
    colframe=ungu,arc=4pt,boxrule=1pt]
  \centering
  File ini \textbf{bukan untuk dikompilasi sendiri.} Ia berisi pengaturan
  bersama (paket, warna, kotak, perintah) yang dipanggil lewat
  \texttt{\textbackslash input} oleh berkas \textbf{MAE-2026-2027-Latihan1.tex}
  sampai \textbf{Latihan9.tex}.

  \medskip
  Silakan buka dan kompilasi salah satu berkas \textbf{Latihan1--Latihan9}
  tersebut untuk menghasilkan PDF pembahasan.
  \end{tcolorbox}
  \end{center}
  \end{document}
\fi

\renewcommand{\babjudul}{Pembahasan Latihan — Bab 7}

\begin{document}
\setcounter{chapter}{7}
\buathalamanjudul{Statistika}{7}

% ============================================================
\chapter*{Pembahasan Bab 7: Statistika}
\addcontentsline{toc}{chapter}{Pembahasan Bab 7: Statistika}
\markboth{PEMBAHASAN BAB 7}{}
% ============================================================

\begin{tips}
\textbf{Cara membaca buku ini.} Tiap soal: \textbf{(1)} soal ditulis ulang,
\textbf{(2)} konsep kunci, \textbf{(3)} langkah demi langkah, \textbf{(4)}
jebakan, \textbf{(5)} variasi. \emph{Aturan emas Bab 7:} \textbf{urutkan data
lebih dulu} sebelum mencari median/kuartil, dan selalu tanya ``apakah rata-rata
saja cukup, atau perlu ukuran sebaran?''.

\smallskip
\textbf{Catatan metode kuartil.} Buku ini memakai \emph{rumus letak}
$\text{letak }Q_i=\dfrac{i(n+1)}{4}$ (dan analog untuk desil/persentil). Jika
letaknya pecahan, lakukan \emph{interpolasi} antara dua datum terdekat.
\end{tips}

% ============================================================
\section*{Studi Kasus: Menganalisis Kebiasaan Belajar Siswa}
\addcontentsline{toc}{section}{Studi Kasus: Kebiasaan Belajar}
% ============================================================

\begin{studikasus}{Menganalisis Kebiasaan Belajar Siswa}
Lama belajar (jam/hari) dua kelompok:
\begin{center}
\begin{tabular}{l|l}
\toprule
Kelompok & Data (jam) \\ \midrule
A & $2,2,3,3,3,3,4,4$ \\
B & $0,1,2,3,3,4,5,6$ \\ \bottomrule
\end{tabular}
\end{center}
\textbf{Pertanyaan.} (a) Hitung rata-rata tiap kelompok --- apakah sama?
(b) Kelompok mana yang ``lebih konsisten''? (c) Ukuran apa yang menangkap
``kekonsistenan'' itu?
\end{studikasus}

\subsection*{Jawaban lengkap studi kasus}

\begin{ingat}{Mean vs sebaran}
\emph{Mean} $\bar x=\frac{\sum x_i}{n}$ hanya menyatakan ``pusat''. Dua data
bisa berpusat sama tapi \emph{beda sebaran}. Ukuran sebaran (jangkauan,
simpangan baku) menangkap ``kekonsistenan''.
\end{ingat}

\textbf{(a) Rata-rata.}
\[
\bar x_A=\frac{2+2+3+3+3+3+4+4}{8}=\frac{24}{8}=3,\qquad
\bar x_B=\frac{0+1+2+3+3+4+5+6}{8}=\frac{24}{8}=3.
\]
Keduanya \textbf{sama}, yaitu $3$ jam/hari.

\textbf{(b) Yang lebih konsisten.} Kelompok \textbf{A} --- nilainya berkerumun
rapat di sekitar $3$, sedangkan B menyebar dari $0$ sampai $6$.

\textbf{(c) Ukuran yang menangkap.} \emph{Simpangan baku} (atau jangkauan).
Hitung untuk masing-masing ($\bar x=3$):
\[
s_A=\sqrt{\frac{1+1+0+0+0+0+1+1}{8}}=\sqrt{\frac{4}{8}}=\sqrt{0{,}5}\approx0{,}71,
\]
\[
s_B=\sqrt{\frac{9+4+1+0+0+1+4+9}{8}}=\sqrt{\frac{28}{8}}=\sqrt{3{,}5}\approx1{,}87.
\]
\jadi meski rata-rata sama, $s_A<s_B$ membuktikan A \textbf{lebih konsisten}
secara kuantitatif. Rata-rata saja tak cukup --- itulah pesan besar bab ini.

% ############################################################
\section*{Subbab 7.1 — Ukuran Pemusatan dan Penyebaran}
\addcontentsline{toc}{section}{Subbab 7.1 — Pemusatan \& Penyebaran}
% ############################################################

\begin{ingat}{Rumus inti}
\[
\bar x=\frac{\sum x_i}{n},\qquad \bar x_{\text{kel}}=\frac{\sum f_ix_i}{\sum f_i},
\qquad s=\sqrt{\frac{\sum(x_i-\bar x)^2}{n}}.
\]
Median = nilai tengah data \emph{terurut}; modus = yang paling sering.
Jangkauan $=x_{\max}-x_{\min}$; JAK $=Q_3-Q_1$. Letak $Q_i=\frac{i(n+1)}{4}$.
\end{ingat}

% ---------------- SOAL 7.1.1 ----------------
\soalno{7.1.1}{kuartil (rumus letak)}
\begin{soalbox}{}
Tentukan $Q_1$, $Q_2$, $Q_3$ dari data $3,5,5,7,8,9,11,12$.
\end{soalbox}

\jawab
\langkah{1} Data sudah terurut, $n=8$. Gunakan letak $Q_i=\dfrac{i(n+1)}{4}
=\dfrac{i\cdot9}{4}$.
\langkah{2} $Q_1$: letak $\dfrac{9}{4}=2{,}25$ $\Rightarrow$ antara datum ke-2
($5$) dan ke-3 ($5$): $Q_1=5+0{,}25(5-5)=5$.
\langkah{3} $Q_2$: letak $\dfrac{18}{4}=4{,}5$ $\Rightarrow$ antara datum ke-4
($7$) dan ke-5 ($8$): $Q_2=7+0{,}5(8-7)=7{,}5$.
\langkah{4} $Q_3$: letak $\dfrac{27}{4}=6{,}75$ $\Rightarrow$ antara datum ke-6
($9$) dan ke-7 ($11$): $Q_3=9+0{,}75(11-9)=9+1{,}5=10{,}5$.
\jadi $Q_1=5,\ Q_2=7{,}5,\ Q_3=10{,}5$.

\begin{jebakan}
\begin{itemize}[leftmargin=*,topsep=2pt]
  \item \textbf{Lupa mengurutkan.} Untuk data acak, urutkan dulu; di sini sudah
  terurut.
  \item \textbf{Salah interpolasi.} ``Letak $6{,}75$'' berarti $75\%$ jarak dari
  datum ke-6 menuju ke-7.
\end{itemize}
\end{jebakan}

\begin{variasi}
Metode ``median dari separuh'' (membelah data) bisa memberi hasil sedikit
berbeda; ikuti metode yang diminta soal. Data ganjil ($n$ ganjil) membuat
median tepat di satu datum.
\end{variasi}

% ---------------- SOAL 7.1.2 ----------------
\soalno{7.1.2}{simpangan baku}
\begin{soalbox}{}
Hitung simpangan baku dari data $2,4,4,6,9$.
\end{soalbox}

\begin{ingat}{Langkah simpangan baku}
(1) Hitung mean. (2) Kurangi mean dari tiap data, kuadratkan. (3) Jumlahkan,
bagi $n$ (variansi). (4) Akarkan.
\end{ingat}

\jawab
\langkah{1} Mean: $\bar x=\dfrac{2+4+4+6+9}{5}=\dfrac{25}{5}=5$.
\langkah{2} Selisih dan kuadratnya:
\[
(2-5)^2=9,\ (4-5)^2=1,\ (4-5)^2=1,\ (6-5)^2=1,\ (9-5)^2=16.
\]
\langkah{3} Jumlah $=9+1+1+1+16=28$; variansi $=\dfrac{28}{5}=5{,}6$.
\langkah{4} $s=\sqrt{5{,}6}\approx2{,}37$.
\jadi $s\approx2{,}37$.

\begin{jebakan}
\begin{itemize}[leftmargin=*,topsep=2pt]
  \item \textbf{Lupa mengkuadratkan} atau lupa mengakarkan di akhir.
  \item \textbf{Membagi $n-1$ vs $n$.} Untuk data \emph{populasi} bagi $n$; untuk
  \emph{sampel} bagi $n-1$. Ikuti konteks (buku ini memakai $n$).
\end{itemize}
\end{jebakan}

\begin{variasi}
Hitung ragam/variansi saja (tanpa akar); atau bandingkan $s$ dua kelompok
(seperti studi kasus).
\end{variasi}

% ---------------- SOAL 7.1.3 ----------------
\soalno{7.1.3}{mean, median, modus}
\begin{soalbox}{}
Tentukan mean, median, dan modus dari $5,7,7,8,10$.
\end{soalbox}

\jawab
\langkah{1} Mean: $\dfrac{5+7+7+8+10}{5}=\dfrac{37}{5}=7{,}4$.
\langkah{2} Median: data terurut $n=5$, nilai tengah = datum ke-3 = $7$.
\langkah{3} Modus: nilai paling sering = $7$ (muncul dua kali).
\jadi mean $=7{,}4$, median $=7$, modus $=7$.

\begin{jebakan}
\begin{itemize}[leftmargin=*,topsep=2pt]
  \item \textbf{Median = rata-rata?} Bukan. Median = nilai \emph{posisi tengah}
  data terurut, bukan hasil bagi.
\end{itemize}
\end{jebakan}

\begin{variasi}
Data tanpa modus (semua beda) atau bimodal (dua modus); data $n$ genap (median
rata-rata dua tengah).
\end{variasi}

% ---------------- SOAL 7.1.4 ----------------
\soalno{7.1.4}{jangkauan \& JAK}
\begin{soalbox}{}
Hitung jangkauan dan jangkauan antarkuartil (JAK) data $3,5,5,7,8,9,11,12$.
\end{soalbox}

\jawab
\langkah{1} Jangkauan $=x_{\max}-x_{\min}=12-3=9$.
\langkah{2} JAK $=Q_3-Q_1$. Dari Soal 7.1.1: $Q_1=5$, $Q_3=10{,}5$.
\langkah{3} JAK $=10{,}5-5=5{,}5$.
\jadi jangkauan $=9$; JAK $=5{,}5$.

\begin{jebakan}
\begin{itemize}[leftmargin=*,topsep=2pt]
  \item \textbf{Menukar jangkauan dan JAK.} Jangkauan pakai nilai ekstrem; JAK
  pakai kuartil (mengukur $50\%$ data tengah, tahan pencilan).
\end{itemize}
\end{jebakan}

\begin{variasi}
Bandingkan jangkauan (peka pencilan) dengan JAK (tahan pencilan) pada data yang
memuat outlier.
\end{variasi}

% ---------------- SOAL 7.1.5 ----------------
\soalno{7.1.5}{mean data berkelompok}
\begin{soalbox}{}
Tentukan mean data berkelompok: nilai $5,6,7,8$ dengan frekuensi $2,4,3,1$.
\end{soalbox}

\begin{ingat}{Mean berbobot frekuensi}
$\bar x=\dfrac{\sum f_i x_i}{\sum f_i}$: tiap nilai ``diberi bobot'' sesuai
seberapa sering muncul.
\end{ingat}

\jawab
\langkah{1} Hitung $\sum f_i x_i$:
\[
5(2)+6(4)+7(3)+8(1)=10+24+21+8=63.
\]
\langkah{2} $\sum f_i=2+4+3+1=10$.
\langkah{3} $\bar x=\dfrac{63}{10}=6{,}3$.
\jadi mean $=6{,}3$.

\begin{jebakan}
\begin{itemize}[leftmargin=*,topsep=2pt]
  \item \textbf{Membagi banyak nilai (4), bukan total frekuensi (10).} Penyebut
  adalah $\sum f_i$, yaitu banyak \emph{datum}, bukan banyak \emph{nilai berbeda}.
  \item \textbf{Lupa mengalikan nilai dengan frekuensinya.}
\end{itemize}
\end{jebakan}

\begin{variasi}
Data berkelompok dengan interval kelas (gunakan titik tengah kelas sebagai
$x_i$); mean gabungan dua kelompok.
\end{variasi}

% ---------------- SOAL 7.1.6 ----------------
\soalno{7.1.6}{mean \& median sederhana}
\begin{soalbox}{}
Data $10,12,12,15,20$. Hitung mean dan median.
\end{soalbox}

\jawab
\langkah{1} Mean: $\dfrac{10+12+12+15+20}{5}=\dfrac{69}{5}=13{,}8$.
\langkah{2} Median: $n=5$ (ganjil), nilai tengah = datum ke-3 = $12$.
\jadi mean $=13{,}8$; median $=12$.

\begin{jebakan}
\begin{itemize}[leftmargin=*,topsep=2pt]
  \item \textbf{Mengira median = rata-rata nilai terkecil dan terbesar.}
  Median = nilai \emph{posisi} tengah.
\end{itemize}
\end{jebakan}

\begin{variasi}
Tambahkan pencilan (mis.\ ganti $20\to200$): mean melonjak, median tetap $12$
--- peragaan ketahanan median.
\end{variasi}

% ---------------- SOAL 7.1.7 ----------------
\soalno{7.1.7}{mean dan penghapusan datum}
\begin{soalbox}{}
Jika rata-rata $5$ bilangan $=8$ dan setelah satu bilangan dibuang rata-rata $4$
sisanya menjadi $7{,}5$, tentukan bilangan yang dibuang.
\end{soalbox}

\begin{ingat}{Mean $\to$ jumlah}
Kunci: ubah rata-rata menjadi \emph{jumlah} ($\text{jumlah}=\bar x\times n$).
Bilangan yang dibuang = selisih dua jumlah.
\end{ingat}

\jawab
\langkah{1} Jumlah $5$ bilangan $=8\times5=40$.
\langkah{2} Jumlah $4$ bilangan sisa $=7{,}5\times4=30$.
\langkah{3} Bilangan yang dibuang $=40-30=10$.
\jadi bilangan yang dibuang $=10$.

\begin{jebakan}
\begin{itemize}[leftmargin=*,topsep=2pt]
  \item \textbf{Mengurangi rata-rata} ($8-7{,}5=0{,}5$). Rata-rata tak bisa
  dikurangi langsung; ubah ke jumlah dulu.
\end{itemize}
\end{jebakan}

\begin{variasi}
Menambah datum baru yang mengubah rata-rata; mencari nilai agar rata-rata
mencapai target tertentu.
\end{variasi}

% ---------------- SOAL 7.1.8 ----------------
\soalno{7.1.8}{modus}
\begin{soalbox}{}
Tentukan modus dari $3,3,4,5,5,5,6$.
\end{soalbox}

\jawab
\langkah{1} Hitung kemunculan: $3$ (2$\times$), $4$ (1$\times$), $5$ (3$\times$),
$6$ (1$\times$).
\langkah{2} Yang terbanyak: $5$.
\jadi modus $=5$.

\begin{jebakan}
\begin{itemize}[leftmargin=*,topsep=2pt]
  \item \textbf{Menganggap modus = nilai tengah.} Modus = nilai yang \emph{paling
  sering}, bukan yang di tengah.
\end{itemize}
\end{jebakan}

\begin{variasi}
Data bimodal (dua nilai sama-sama terbanyak); data tanpa modus.
\end{variasi}

% ---------------- SOAL 7.1.9 ----------------
\soalno{7.1.9}{median data ganjil}
\begin{soalbox}{}
Hitung $Q_2$ dari $1,3,3,6,7,8,9$.
\end{soalbox}

\jawab
\langkah{1} $Q_2=$ median. Data terurut, $n=7$. Letak $Q_2=\dfrac{2(7+1)}{4}
=\dfrac{16}{4}=4$ $\Rightarrow$ datum ke-4.
\langkah{2} Datum ke-4 $=6$.
\jadi $Q_2=6$.

\begin{jebakan}
\begin{itemize}[leftmargin=*,topsep=2pt]
  \item \textbf{Mengambil datum ke-$\frac n2$.} Untuk $n$ ganjil, posisi tengah
  adalah $\frac{n+1}{2}=4$ (datum ke-4), bukan ke-3,5.
\end{itemize}
\end{jebakan}

\begin{variasi}
$n$ genap $\Rightarrow$ median rata-rata dua datum tengah.
\end{variasi}

% ---------------- SOAL 7.1.10 ----------------
\soalno{7.1.10}{data simetris}
\begin{soalbox}{}
Sebuah data simetris: apa hubungan mean, median, dan modusnya?
\end{soalbox}

\jawab
\langkah{1} Pada distribusi \emph{simetris} (dan unimodal), pusatnya berimpit.
\jadi \textbf{mean $=$ median $=$ modus}. (Sebaliknya, jika ketiganya berbeda,
data \emph{miring/skewed}: pada miring kanan mean $>$ median $>$ modus.)

\begin{center}
\begin{tikzpicture}[scale=0.9]
\begin{axis}[axis lines=middle,width=8cm,height=4cm,xtick=\empty,ytick=\empty,
  xlabel=$x$,ymin=0,ymax=1.1,domain=-3:3,samples=60,clip=false]
\addplot[very thick,biru]{exp(-x^2/1.2)};
\draw[dashed,merah,thick] (axis cs:0,0)--(axis cs:0,1);
\node[merah,font=\scriptsize,above] at (axis cs:0,1){mean=median=modus};
\end{axis}
\end{tikzpicture}
\end{center}

\begin{jebakan}
\begin{itemize}[leftmargin=*,topsep=2pt]
  \item \textbf{Mengira selalu sama untuk data apa pun.} Kesamaan ini khusus
  distribusi \emph{simetris}; data miring memisahkan ketiganya.
\end{itemize}
\end{jebakan}

\begin{variasi}
Tentukan arah kemiringan dari posisi relatif mean vs median; sketsa kurva miring.
\end{variasi}

% ############################################################
\section*{Subbab 7.2 — Statistika Lanjutan (Letak, Box Plot, Outlier, Z-score)}
\addcontentsline{toc}{section}{Subbab 7.2 — Statistika Lanjutan}
% ############################################################

\begin{ingat}{Ukuran letak, outlier, z-score}
Letak $Q_i=\frac{i(n+1)}{4}$, $D_i=\frac{i(n+1)}{10}$, $P_i=\frac{i(n+1)}{100}$
(interpolasi bila pecahan). Hubungan: $Q_2=D_5=P_{50}=$ median; $Q_1=P_{25}$,
$Q_3=P_{75}$.\\
\textbf{Outlier:} di luar pagar $Q_1-1{,}5\,\text{JAK}$ atau $Q_3+1{,}5\,
\text{JAK}$.\quad \textbf{Z-score:} $z=\dfrac{x-\bar x}{s}$.
\end{ingat}

% ---------------- SOAL 7.2.1 ----------------
\soalno{7.2.1}{kuartil data ganjil}
\begin{soalbox}{}
Tentukan $Q_1,Q_2,Q_3$ dari $2,4,5,7,8,10,12$.
\end{soalbox}

\jawab
\langkah{1} Terurut, $n=7$, $n+1=8$. Letak $Q_i=\dfrac{i\cdot8}{4}=2i$.
\langkah{2} $Q_1$: letak $2$ $\Rightarrow$ datum ke-2 $=4$.
\langkah{3} $Q_2$: letak $4$ $\Rightarrow$ datum ke-4 $=7$.
\langkah{4} $Q_3$: letak $6$ $\Rightarrow$ datum ke-6 $=10$.
\jadi $Q_1=4,\ Q_2=7,\ Q_3=10$ (semua tepat di datum karena $n+1=8$ habis dibagi
$4$).

\begin{jebakan}
\begin{itemize}[leftmargin=*,topsep=2pt]
  \item \textbf{Salah hitung letak.} $\frac{i(n+1)}{4}$ memakai $n+1=8$, bukan
  $n=7$.
\end{itemize}
\end{jebakan}

\begin{variasi}
$n$ yang membuat letak pecahan (perlu interpolasi), seperti Soal 7.1.1.
\end{variasi}

% ---------------- SOAL 7.2.2 ----------------
\soalno{7.2.2}{JAK / IQR}
\begin{soalbox}{}
Hitung JAK (IQR) data $2,4,5,7,8,10,12$.
\end{soalbox}

\jawab
\langkah{1} Dari Soal 7.2.1: $Q_1=4$, $Q_3=10$.
\langkah{2} JAK $=Q_3-Q_1=10-4=6$.
\jadi JAK $=6$.

\begin{jebakan}
\begin{itemize}[leftmargin=*,topsep=2pt]
  \item \textbf{Memakai min--maks.} JAK khusus $Q_3-Q_1$ (bukan jangkauan
  $x_{\max}-x_{\min}=10$).
\end{itemize}
\end{jebakan}

\begin{variasi}
JAK sebagai dasar deteksi outlier (Soal 7.2.6) dan lebar kotak box plot.
\end{variasi}

% ---------------- SOAL 7.2.3 ----------------
\soalno{7.2.3}{desil (interpolasi)}
\begin{soalbox}{}
Tentukan $D_4$ dari $1,3,4,6,7,9,10,12,14,15$.
\end{soalbox}

\jawab
\langkah{1} $n=10$, $n+1=11$. Letak $D_4=\dfrac{4(n+1)}{10}=\dfrac{4\cdot11}{10}
=4{,}4$.
\langkah{2} Letak $4{,}4$ $\Rightarrow$ antara datum ke-4 ($6$) dan ke-5 ($7$),
sejauh $0{,}4$ dari datum ke-4:
\[
D_4=6+0{,}4(7-6)=6+0{,}4=6{,}4.
\]
\jadi $D_4=6{,}4$.

\begin{jebakan}
\begin{itemize}[leftmargin=*,topsep=2pt]
  \item \textbf{Membulatkan letak ke datum ke-4 saja} ($=6$) tanpa interpolasi.
  Bagian pecahan ($0{,}4$) penting.
\end{itemize}
\end{jebakan}

\begin{variasi}
$D_5$ (= median); persentil (Soal 7.2.4). Semua memakai kerangka letak yang sama.
\end{variasi}

% ---------------- SOAL 7.2.4 ----------------
\soalno{7.2.4}{persentil (interpolasi)}
\begin{soalbox}{}
Tentukan $P_{25}$ dari data $5,8,10,12,15,18,20,24$.
\end{soalbox}

\jawab
\langkah{1} $n=8$, $n+1=9$. Letak $P_{25}=\dfrac{25(n+1)}{100}=\dfrac{25\cdot9}{100}
=2{,}25$.
\langkah{2} Antara datum ke-2 ($8$) dan ke-3 ($10$), sejauh $0{,}25$:
\[
P_{25}=8+0{,}25(10-8)=8+0{,}5=8{,}5.
\]
\jadi $P_{25}=8{,}5$ (yang juga sama dengan $Q_1$, karena $P_{25}=Q_1$).

\begin{jebakan}
\begin{itemize}[leftmargin=*,topsep=2pt]
  \item \textbf{Menghitung $25\%$ dari nilai}, bukan letak. $P_{25}$ menunjuk
  \emph{posisi}, lalu ambil nilainya lewat interpolasi.
\end{itemize}
\end{jebakan}

\begin{variasi}
$P_{50}$ (median), $P_{75}$ ($=Q_3$); persentil dari nilai ujian sekelas.
\end{variasi}

% ---------------- SOAL 7.2.5 ----------------
\soalno{7.2.5}{lima serangkai}
\begin{soalbox}{}
Sebutkan lima serangkai dari $3,6,7,9,12,15,20$.
\end{soalbox}

\begin{ingat}{Lima serangkai}
Min, $Q_1$, $Q_2$ (median), $Q_3$, Maks --- ringkasan yang mendasari box plot.
\end{ingat}

\jawab
\langkah{1} $n=7$, $n+1=8$, letak $Q_i=2i$.
\langkah{2} Min $=3$; $Q_1=$ datum ke-2 $=6$; $Q_2=$ datum ke-4 $=9$;
$Q_3=$ datum ke-6 $=15$; Maks $=20$.
\jadi lima serangkainya: $(3,\ 6,\ 9,\ 15,\ 20)$.

\begin{jebakan}
\begin{itemize}[leftmargin=*,topsep=2pt]
  \item \textbf{Melupakan min/maks.} Lima serangkai mencakup \emph{dua} nilai
  ekstrem, bukan hanya kuartil.
\end{itemize}
\end{jebakan}

\begin{variasi}
Gunakan lima serangkai untuk menggambar box plot (Soal 7.2.10).
\end{variasi}

% ---------------- SOAL 7.2.6 ----------------
\soalno{7.2.6}{deteksi outlier (aturan pagar)}
\begin{soalbox}{}
Data $10,11,12,13,40$ dengan $Q_1=11$, $Q_3=13$. Apakah $40$ pencilan?
\end{soalbox}

\jawab
\langkah{1} JAK $=Q_3-Q_1=13-11=2$.
\langkah{2} Pagar atas $=Q_3+1{,}5\cdot\text{JAK}=13+1{,}5(2)=13+3=16$.
\langkah{3} Bandingkan: $40>16$.
\jadi \textbf{ya}, $40$ adalah pencilan (di luar pagar atas).

\begin{jebakan}
\begin{itemize}[leftmargin=*,topsep=2pt]
  \item \textbf{Menebak dari ``kelihatan besar''.} Gunakan aturan pagar objektif
  ($1{,}5\times$ JAK), bukan perasaan.
\end{itemize}
\end{jebakan}

\begin{variasi}
Cek pagar bawah $Q_1-1{,}5\,\text{JAK}=11-3=8$; nilai $<8$ juga pencilan.
\end{variasi}

% ---------------- SOAL 7.2.7 ----------------
\soalno{7.2.7}{z-score}
\begin{soalbox}{}
Hitung $z$-score nilai $90$ jika $\bar x=75$ dan $s=10$.
\end{soalbox}

\begin{ingat}{Z-score = jarak dalam satuan simpangan baku}
$z=\dfrac{x-\bar x}{s}$: ``berapa simpangan baku'' nilai dari rata-rata.
$z>0$ di atas rata-rata; $z<0$ di bawah.
\end{ingat}

\jawab
\langkah{1} $z=\dfrac{x-\bar x}{s}=\dfrac{90-75}{10}=\dfrac{15}{10}=1{,}5$.
\jadi $z=1{,}5$ (nilai $90$ berada $1{,}5$ simpangan baku \emph{di atas}
rata-rata).

\begin{jebakan}
\begin{itemize}[leftmargin=*,topsep=2pt]
  \item \textbf{Membagi $\bar x$ bukan $s$.} Penyebutnya simpangan baku $s$.
\end{itemize}
\end{jebakan}

\begin{variasi}
Membandingkan dua mapel berskala beda lewat z-score; nilai dengan $z$ negatif.
\end{variasi}

% ---------------- SOAL 7.2.8 ----------------
\soalno{7.2.8}{z-score terbalik (cari nilai asli)}
\begin{soalbox}{}
Nilai $z=-1{,}5$ pada data $\bar x=60$, $s=8$. Tentukan nilai aslinya.
\end{soalbox}

\jawab
\langkah{1} Balik rumus: dari $z=\dfrac{x-\bar x}{s}$ diperoleh
$x=\bar x+z\cdot s$.
\langkah{2} $x=60+(-1{,}5)(8)=60-12=48$.
\jadi nilai aslinya $48$.

\begin{jebakan}
\begin{itemize}[leftmargin=*,topsep=2pt]
  \item \textbf{Salah tanda.} $z$ negatif $\Rightarrow$ nilai \emph{di bawah}
  rata-rata ($48<60$).
\end{itemize}
\end{jebakan}

\begin{variasi}
Diberi $z$ positif; atau mencari $s$/$\bar x$ dari $z$ dan $x$ yang diketahui.
\end{variasi}

% ---------------- SOAL 7.2.9 ----------------
\soalno{7.2.9}{makna persentil}
\begin{soalbox}{}
Seseorang berada di persentil ke-$90$ tinggi badan. Jelaskan maknanya.
\end{soalbox}

\jawab
\langkah{1} Persentil ke-$90$ ($P_{90}$) berarti posisi yang $90\%$ data berada
\emph{di bawahnya}.
\jadi orang itu \textbf{lebih tinggi daripada sekitar $90\%$ orang} dalam
kelompoknya; hanya sekitar $10\%$ yang lebih tinggi. (Ini ukuran \emph{posisi
relatif}, bukan nilai tinggi mutlak.)

\begin{jebakan}
\begin{itemize}[leftmargin=*,topsep=2pt]
  \item \textbf{Mengira ``$90\%$ dari tinggi maksimum''.} Persentil mengukur
  peringkat/posisi, bukan persentase nilai.
\end{itemize}
\end{jebakan}

\begin{variasi}
Persentil pertumbuhan anak (grafik KMS); persentil skor ujian nasional.
\end{variasi}

% ---------------- SOAL 7.2.10 ----------------
\soalno{7.2.10}{menggambar box plot}
\begin{soalbox}{}
Gambarkan sketsa box plot untuk data $2,3,5,5,6,8,9,14$ (tentukan dulu lima
serangkainya).
\end{soalbox}

\jawab
\langkah{1} Terurut, $n=8$, $n+1=9$, letak $Q_i=\dfrac{9i}{4}$.
\langkah{2} $Q_1$: letak $2{,}25$ $\Rightarrow$ antara datum ke-2 ($3$) dan
ke-3 ($5$): $Q_1=3+0{,}25(2)=3{,}5$.
$Q_2$: letak $4{,}5$ $\Rightarrow$ antara datum ke-4 ($5$) dan ke-5 ($6$):
$Q_2=5{,}5$.
$Q_3$: letak $6{,}75$ $\Rightarrow$ antara datum ke-6 ($8$) dan ke-7 ($9$):
$Q_3=8+0{,}75(1)=8{,}75$.
\langkah{3} Min $=2$, Maks $=14$. Lima serangkai: $(2;\ 3{,}5;\ 5{,}5;\ 8{,}75;\
14)$.

\begin{center}
\begin{tikzpicture}[scale=0.7]
  \draw[->](0,0)--(15.5,0) node[right]{nilai};
  \foreach \x in {0,2,...,14}{\draw (\x,-0.08)--(\x,0.08) node[below,font=\tiny]{\x};}
  \draw[very thick,biru] (3.5,0.5) rectangle (8.75,1.5);
  \draw[very thick,biru] (5.5,0.5)--(5.5,1.5);
  \draw[very thick,biru] (2,1)--(3.5,1);
  \draw[very thick,biru] (8.75,1)--(14,1);
  \node[above,font=\scriptsize] at (2,1){min $2$};
  \node[above,font=\scriptsize] at (3.5,1.55){$Q_1{=}3{,}5$};
  \node[above,font=\scriptsize] at (5.5,1.55){$Q_2{=}5{,}5$};
  \node[above,font=\scriptsize] at (8.75,1.55){$Q_3{=}8{,}75$};
  \node[above,font=\scriptsize] at (14,1){maks $14$};
\end{tikzpicture}
\end{center}
\jadi box plot memiliki kotak dari $3{,}5$ ke $8{,}75$ (median $5{,}5$), dengan
kumis menjangkau $2$ dan $14$.

\begin{jebakan}
\begin{itemize}[leftmargin=*,topsep=2pt]
  \item \textbf{Menggambar kotak dari min ke maks.} Kotak hanya $Q_1$--$Q_3$
  ($50\%$ data tengah); ekstrem digambar sebagai kumis.
\end{itemize}
\end{jebakan}

\begin{variasi}
Data dengan pencilan (titik terpisah di luar kumis); membandingkan dua box plot
berdampingan.
\end{variasi}

% ============================================================
\begin{tips}
\textbf{Ringkasan strategi Bab 7.}
\begin{enumerate}[leftmargin=*,topsep=2pt]
  \item \emph{Urutkan dulu} sebelum median/kuartil. Kuartil/desil/persentil pakai
  \textbf{rumus letak} $\frac{i(n+1)}{k}$ + interpolasi.
  \item \emph{Mean} peka pencilan; \emph{median} tahan pencilan. Pilih sesuai
  konteks (gaji direktur $\to$ median).
  \item \emph{Sebaran:} jangkauan (peka), JAK $=Q_3-Q_1$ (tahan), simpangan baku
  (kuadratkan selisih agar tak saling menghapus, lalu akarkan).
  \item \emph{Outlier} lewat pagar $1{,}5\times$JAK; \emph{z-score} untuk
  membandingkan skala berbeda; \emph{box plot} merangkum lima serangkai.
\end{enumerate}
\end{tips}

\end{document}
